\mathplain
\section{Designs That Did Not Work}\label{sec:history}

Several parts of \FCdot look arbitrary until one sees what they replaced. This section records the designs that were tried first and why each was abandoned.

\subsection{Existential packages}\label{sec:counterexample}

The first design followed F-ing modules \cite{rossberg2014fing}, which interprets ML modules by existential packages in System $F_\omega$. A declaration $\dtyp{A}{S}{T}$ would become an existential package, and opening a variable $x$ of that type would yield a fresh type name $\alpha$ with $\src{S} \le \alpha \le \src{T}$. Every later fact about $\dsel{x}{A}$ would be a fact about $\alpha$. Example~\ref{ex:mini-cast} shows why this does not work. In the context
\[
\src{S = \dtyp{A}{\bot}{\top}, \qquad T = \dtyp{A}{\mathit{Int}}{\top}, \qquad \Gamma = x : \dtyp{B}{S}{T},\ w : S}
\]
DOT derives $\src{\mathit{Int} \sub w.A}$ through $\src{w : T}$, so
\[
\src{\klet g = \lambda(y : w.A).\,y \kin g\,n} \qquad \text{with } \src{n : \mathit{Int}}
\]
is well typed. In the package design, $\src{w.A}$ denotes the name $\alpha$ opened from the declared type $\src{S}$ of $\src{w}$. The retyping of $\src{w}$ at $\src{T}$ is a coercion between two package types, and using it means opening the coerced package, which yields a fresh $\alpha'$ with $\src{\mathit{Int}} \le \alpha'$. Nothing relates $\alpha'$ to $\alpha$, and the application $\src{g\,n}$ has no translation. The context is unrealizable, but a total translation on derivations has to translate this derivation anyway.

The conclusion was that the identity of a member must be attached to the binder, not to the type through which the binder is viewed. That is what block names do. In \FCdot the type $\src{w.A}$ is the name $\sel{w}{A}$ whatever type $w$ is viewed at, and the fact $\mathit{Int} \le \sel{w}{A}$ is an elimination at the cast atom $\ev{cast}\,w\,e_{ST}$, whose root is $w$. This program, example E4 of the development, became the acceptance test of every later design. Whether the sharing machinery of F-ing modules can be adapted to lower bounds we have not investigated.

\subsection{Morphisms that eliminate through the self}\label{sec:why-structural}

The second design let an object coercion prove a target proposition by arbitrary evidence under the self binder:
\[
\tgt{\Gamma, z : \obj{\tel} \vdash k : S \le T} .
\]
This is more expressive than templates. A morphism body $k$ may eliminate at $z$, compose the results, and eliminate again. Its normal forms are closures. To use such a morphism at an atom $a$, the normalizer has to substitute, $k[z \mapsto a]$, and normalize the result. But $k[z \mapsto a]$ is not a subterm of anything, and the normalization theorem has no induction hypothesis for it. It may perform further eliminations through the atom, whose views come from further morphisms. Composition of two such morphisms needs an environment of views, and the typedness of a form needs a step index with an account of the eliminations performed.

We proved a good part of that development before observing that the DOT translation never produces such a morphism. Every morphism it produces reads a source proposition and adjusts it by closed evidence. Restricting morphisms to templates loses nothing for DOT and turns the normalization theorem into an induction on syntax: the entries of an object form do not depend on the atom they are applied to, composition substitutes entries into entries, and application to an atom looks up named propositions in a view. The two places where templates had to grow later, composition templates and bound entries, keep that property. A composition template chains two templates over the same source, and a bound entry is closed evidence out of the source as a whole.

\subsection{Well-founded resolution and the self-alias restriction}

The first version of resolution followed definitions by a well-founded measure. That required that no witness of a literal be a bare name of the same block, a side condition on the literal rule called the self-alias restriction. The condition propagated into the translation, where a field declared at $\dsel{x}{A}$ inside its own literal makes the witness of $\sel{x}{a}$ the name $\sel{x}{A}$, and from there into a restriction on the source rule \rn{\{\}-I}. It rejected literals such as
\[
\src{\nu(x.\,\ddeftyp{T}{\mathit{Int}} \wedge \ddeftrm{v}{n})} \quad\text{typed at}\quad \src{\dmu{x}{\dtyp{T}{\mathit{Int}}{\mathit{Int}} \wedge \dfld{v}{x.T}}} ,
\]
which DOT accepts. The resolution of Section~\ref{sec:resolution} runs with fuel instead, and a pigeonhole argument shows that the fuel is enough. It accepts every literal, including a two-element cycle, and the restriction disappeared from the target's literal rule, from the source's \rn{\{\}-I}, and from the translation.

\subsection{Restricted shapes}

The first version of the translation restricted the operands of intersections and the bodies of recursive types to declarations, because the target had no telescope for anything else. An intersection with a selection as an operand had no translation.

Self-bounds lifted the restriction on intersections. What they cost on the target side is the routed entries of Section~\ref{sec:canonical}: the pairing of a bound cast with another coercion has entries that prove propositions of an object type other than the source, which an $\form{obj}$ form cannot express.

The bodies of recursive types stayed restricted longer, because a body that is not a declaration becomes a bound that mentions its own self, and \FCdot then had no rule that proves or uses such a bound. Lifting that restriction took bounds that may mention the self, with a cast through a bound reading only a closed one, and the adapters $\mathit{into}$ and $\mathit{out}$ in the translation of \rn{Rec-I} and \rn{Rec-E}. Preservation and progress needed no new case for it. The clause of $\mathit{tel}$ for a recursive type still asks whether the body is a declaration, but only to choose between the telescope of the body and a single closed bound, so that every bound $\mathit{tel}$ produces is closed, which the identity morphisms of \rn{And}$_1$ and \rn{And}$_2$ need.

\subsection{Smaller things the proofs forced}

Two details of the translation were found while proving it, not while designing it. The typedness theorem of the translation needs well-formed contexts, in which the declaration type of every literal binder has exact type members and distinct labels, because the typedness of $\tgt{\mathit{litCo}}$ reads witnesses off the declaration type by position. And the fields of a translated intersection are ordered with the right conjunct outermost, because the right conjunct shadows in \DOTMNF, and both lookup in the target and erasure into the runtime prefer the outermost field.

\subsection{Four revisions of the target}

The target was revised four times while the translation was being proven: to restrict morphisms to templates, to make resolution tolerate aliases, to add self-bounds so that intersections need no restriction, and to let bounds mention the self so that recursive types need none. None of the four revisions added a case to preservation. The third added two cases to an inversion lemma of progress, both impossible. The work of each went into the rules of the target, the normalizer and the lemmas about forms, and the translation. This is the pattern the report has argued for: the work of a new feature goes into evidence and its normal forms, and preservation and progress only ask the canonical forms theorem for a shape.
